\scanpage{206}
\mspage{80}
ens. petit, à cause des hyp. faites sur $M$), les ordinaux $\alpha$ pour lesquels $A_\alpha$ existe satisfont $\operatorname{card}A_\alpha\leq\unsure{\pi|R}$. Donc il doit exister un \deleted{$\unsure{A_{\alpha}}$} $\alpha$ tel que $A_\alpha=B$. Mais alors $i_\alpha\unsure{=|:|\doteq} A\longrightarrow A_\alpha$ est égal à $f$, donc $f\in\widetilde{C^\pi}$, qed

Il reste donc à prouver 6.4. Pour ceci, soit $I$ l'ens des sous-objets de $B$ qui sont dans $\mathrm{Filt}^\pi$ \deleted{$\unsure{=C^\pi}$}, ordonné par inclusion. (Si $B$ est dans $\mathrm{Filt}^\pi$, $A$ aussi, donc $f$ est \inserted{déjà} dans $C^\pi$\inserted{$\subset\unsure{C^\pi|\widetilde{C^\pi}}$}, et il n'y a rien à prouver.) C'est un ens. filtrant grand devant $\pi$, \inserted{(et $=$ $\pi$-\uline{adapté})} dont la $\underrightarrow{\lim}$ est $B$ (5.3. p. \unsure{45|46}). Considérons les
\[
 A_i=B_i\cap A
\]
ils sont dans $\mathrm{Filt}^\pi$, \deleted{\unsure{comme}} sous-objets d'objets de $\mathrm{Filt}^\pi$ (condition d) sur \unsure{la} filtration cardinale). Ainsi, la flèche $f$ apparaît comme limite inductive filtrante des flèches \deleted{$\unsure{A_i\to f_i}$}
\[
 f_i:A_i\to B_i .
\]
Il est \unsure{triviale|trivial} que le système des $A_i$ est cofinal dans $\mathrm{Ss}_\pi(A)$, donc
\[
 A=\underrightarrow{\lim}A_i,\quad\text{d'où}\quad f=\underrightarrow{\lim}f_i.
\]
% archive sheet number: 205
